% ==========================================================================
% COURS 01 -- EXERCICES
% ==========================================================================

\section{Exercices du chapitre}\label{sec:is-exercices}

\begin{TLExercise}{Exercice 1 -- Repêchage d'un véhicule}
Un camion-grue doit relever une voiture tombée dans un port. La première tentative provoque le basculement du camion ; un second véhicule équipé de stabilisateurs est nécessaire.
\begin{enumerate}
  \item Formuler le besoin sans imposer de solution technique.
  \item Définir la frontière du système étudié.
  \item Identifier au moins six éléments de l'environnement.
  \item Proposer quatre exigences mesurables liées à la sécurité et à la performance.
  \item Distinguer les causes relevant d'un composant et celles relevant de l'intégration du système.
\end{enumerate}
\end{TLExercise}

\begin{TLExercise}{Exercice 2 -- Compatibilité d'un train avec les quais}
Une nouvelle rame est plus large afin d'améliorer le confort, mais certaines infrastructures anciennes ne respectent pas un gabarit uniforme.
\begin{enumerate}
  \item Identifier les parties prenantes et les interfaces critiques.
  \item Écrire une exigence de compatibilité vérifiable.
  \item Proposer une méthode de vérification avant la fabrication en série.
  \item Expliquer la différence entre vérification de la rame et validation du service de transport.
\end{enumerate}
\end{TLExercise}

\begin{TLExercise}{Exercice 3 -- Drone de loisir}
Un drone doit permettre à un utilisateur de filmer une scène en extérieur.
\begin{enumerate}
  \item Formuler la mission du système.
  \item Construire la liste des acteurs d'un diagramme de contexte.
  \item Proposer cinq cas d'utilisation formulés par un verbe et un complément.
  \item Écrire six exigences dont deux contraintes réglementaires ou environnementales.
  \item Associer à chaque exigence une méthode de preuve.
\end{enumerate}
\end{TLExercise}

\begin{TLExercise}{Exercice 4 -- Micro-ondes}
On étudie un four à micro-ondes domestique pendant sa phase d'utilisation.
\begin{enumerate}
  \item Définir la matière d'œuvre entrante et sortante.
  \item Énoncer la fonction globale.
  \item Identifier les flux de matière, d'énergie et d'information.
  \item Répartir les fonctions dans les chaînes d'information et d'énergie.
  \item Citer trois autres phases de vie nécessitant un contexte différent.
\end{enumerate}
\end{TLExercise}

\begin{TLExercise}{Exercice 5 -- Diagramme d'exigences}
La mission d'une plateforme élévatrice est de transporter une personne entre deux niveaux. Les exigences initiales sont : charge minimale de 120 kg, vitesse inférieure à 0,15 m/s, arrêt d'urgence en moins de 0,4 s et niveau sonore inférieur à 65 dB(A).
\begin{enumerate}
  \item Attribuer un identifiant à chaque exigence.
  \item Distinguer exigences fonctionnelles et contraintes.
  \item Proposer une exigence globale dont elles dérivent.
  \item Associer à chacune un critère, un niveau, une tolérance et une méthode de vérification.
  \item Représenter les relations «deriveReqt», «satisfy» et «verify».
\end{enumerate}
\end{TLExercise}

\begin{TLExercise}{Exercice 6 -- BDD et IBD d'un axe motorisé}
Un axe comprend un contrôleur, un variateur, un moteur, un réducteur, un chariot et un capteur de position.
\begin{enumerate}
  \item Proposer la décomposition du système dans un BDD.
  \item Identifier les relations de composition.
  \item Décrire les flux échangés dans un IBD.
  \item Préciser la nature et l'unité de chaque flux.
  \item Localiser les interfaces qui devraient faire l'objet d'une spécification détaillée.
\end{enumerate}
\end{TLExercise}

\begin{TLExercise}{Exercice 7 -- Scénario de démarrage}
Un robot mobile reçoit un ordre de départ, vérifie la fermeture de son capot, initialise ses capteurs, libère son frein puis commence à se déplacer. En cas de défaut, il signale une alarme et reste immobile.
\begin{enumerate}
  \item Représenter le scénario sous forme de diagramme de séquence.
  \item Proposer un diagramme d'activités avec décision et branche de défaut.
  \item Définir les états principaux d'un diagramme d'états-transitions.
  \item Écrire deux transitions sous la forme événement [garde] / effet.
\end{enumerate}
\end{TLExercise}

\begin{TLExercise}{Exercice 8 -- Système souhaité, simulé et réel}
La simulation d'un vérin prévoit une course de 500 mm en 4,5 s. Le cahier des charges impose une course d'au moins 490 mm en moins de 5 s. Les mesures donnent 486 mm en 4,8 s.
\begin{enumerate}
  \item Identifier les données relevant des systèmes souhaité, simulé et réel.
  \item Quantifier les trois écarts pertinents.
  \item Conclure sur chaque exigence.
  \item Proposer deux causes possibles de l'écart simulation-réel.
  \item Indiquer si le modèle, le produit ou l'exigence doit être révisé.
\end{enumerate}
\end{TLExercise}

\begin{TLExercise}{Exercice 9 -- Cycle de vie et éco-conception}
Deux solutions de carter sont comparées : aluminium recyclable plus lourd, ou composite plus léger mais difficile à séparer en fin de vie.
\begin{enumerate}
  \item Lister les étapes du cycle de vie à comparer.
  \item Identifier les flux et impacts principaux de chaque solution.
  \item Expliquer pourquoi la seule masse en phase d'usage ne suffit pas pour conclure.
  \item Proposer des indicateurs environnementaux et économiques.
  \item Définir les données supplémentaires nécessaires à la décision.
\end{enumerate}
\end{TLExercise}

\begin{TLExercise}{Exercice 10 -- Modification d'une batterie}
La capacité de la batterie d'un robot est augmentée de 40\%. La nouvelle batterie est plus lourde, plus volumineuse et délivre une tension nominale différente.
\begin{enumerate}
  \item Identifier les exigences susceptibles d'être affectées.
  \item Recenser les interfaces à vérifier.
  \item Proposer une analyse des impacts sur structure, motorisation, autonomie, sécurité et logiciel.
  \item Lister les documents et modèles à mettre à jour.
  \item Définir les essais de non-régression nécessaires.
\end{enumerate}
\end{TLExercise}

\begin{TLExercise}{Exercice 11 -- Architecture fonctionnelle d'un portail automatique}
Un portail doit détecter une demande, vérifier l'absence d'obstacle, déverrouiller, déplacer le vantail, ralentir en fin de course et informer l'utilisateur.
\begin{enumerate}
  \item Formuler la fonction globale.
  \item Construire une décomposition fonctionnelle de type FAST.
  \item Répartir les fonctions dans les chaînes d'information et d'énergie.
  \item Associer à chaque fonction une ou plusieurs solutions techniques possibles.
  \item Identifier les boucles d'information liées à la sécurité.
\end{enumerate}
\end{TLExercise}

\begin{TLExercise}{Exercice 12 -- Étude de synthèse : véhicule autonome}
Un véhicule autonome urbain doit transporter quatre personnes sur un trajet prédéfini, respecter le code de la route et fonctionner par tous les temps raisonnables.
\begin{enumerate}
  \item Identifier les principales parties prenantes et construire le contexte d'utilisation.
  \item Proposer les cas d'utilisation principaux et secondaires.
  \item Rédiger dix exigences mesurables couvrant performance, sécurité, environnement et maintenance.
  \item Définir une architecture fonctionnelle et les flux principaux.
  \item Proposer une décomposition en sous-systèmes pour un BDD.
  \item Décrire un scénario nominal et un scénario de défaillance.
  \item Construire une matrice reliant exigences, sous-systèmes et méthodes de preuve.
  \item Expliquer comment les systèmes souhaité, simulé et réel seront comparés.
\end{enumerate}
\end{TLExercise}
